\documentclass[11pt,a4paper]{ctexart}

\usepackage{amsmath,amssymb}
\usepackage{geometry}
\usepackage{enumitem}
\usepackage{tikz}
\usepackage{listings}
\usepackage{xcolor}
\usepackage{booktabs}
\usepackage{array}
\usepackage{longtable}
\usepackage{tabularx}
\usepackage{multirow}
\usepackage{graphicx}

\geometry{left=2.2cm,right=2.2cm,top=2cm,bottom=2cm}

\lstset{
  basicstyle=\ttfamily\small,
  frame=single,
  language=C,
  showstringspaces=false,
  breaklines=true,
  columns=fullflexible,
}

\setlist[enumerate]{itemsep=0.35em,topsep=0.4em,leftmargin=1.6em}
\setlist[enumerate,2]{label=（\arabic*）,leftmargin=1.8em}

\newcommand{\blank}{\underline{\hspace{3.2em}}}
\newcommand{\blanklong}{\underline{\hspace{6em}}}

\title{}
\author{}
\date{}

\begin{document}

\begin{center}
  {\LARGE\bfseries 哈尔滨工业大学（深圳）}\\[0.6em]
  {\Large\bfseries 2026 年春《计算机系统》期末考试试题（A）}\\[0.3em]
  {\large 回忆还原版}
\end{center}

\vspace{0.4em}
\noindent\textbf{说明：}
\begin{enumerate}[leftmargin=1.5em,itemsep=0.15em,topsep=0.2em]
  \item 本题卷根据考后回忆与课程作业原题整理，非官方原卷。试卷中的具体数值、选项表述、图表数据可能与原卷有一定出入，仅供复习参考。
  \item 选择题另有一题回忆不起来。
\end{enumerate}

\noindent\textbf{考生须知：}本次考试为闭卷考试，考试时间为 120 分钟。

\vspace{0.6em}
\section*{一、选择题}

\begin{enumerate}

\item 关于数据单位，下列说法正确的是（\hspace{1.2em}）。
\begin{enumerate}[label=\Alph*.,itemsep=0.1em]
  \item 1b 表示 1 个字节（Byte）
  \item 1B 表示 1 个比特（bit）
  \item word（字）的位数在所有体系结构中固定为 32 位
  \item word（字）的位数随体系结构而不同（如 32 位机与 64 位机的字长可以不同）
\end{enumerate}

\item 补码加法运算的溢出判别中，以下说法正确的是（\hspace{1.2em}）。
\begin{enumerate}[label=\Alph*.,itemsep=0.1em]
  \item 符号相同的两个数相加必定不会发生溢出
  \item 符号不同的两个数相加可能发生溢出
  \item 符号相同的两个数相加必定发生溢出
  \item 符号不同的两个数相加不可能发生溢出
\end{enumerate}

\item 关于递归函数调用与栈帧，下列说法错误的是（\hspace{1.2em}）。
\begin{enumerate}[label=\Alph*.,itemsep=0.1em]
  \item 每次递归调用通常会在栈上分配一个新的栈帧
  \item 递归深度过大可能导致栈溢出
  \item 同一函数的不同递归调用共享同一个栈帧中的局部变量存储
  \item 返回地址一般保存在对应栈帧中，用于函数返回
\end{enumerate}

\item Y86-64 的 PIPE 流水线实现中，流水线划分为（\hspace{1.2em}）个阶段。
\begin{enumerate}[label=\Alph*.,itemsep=0.1em]
  \item 3
  \item 4
  \item 5
  \item 6
\end{enumerate}

\item 在存储器层次结构中，下列不属于 CPU 片上存储的是（\hspace{1.2em}）。
\begin{enumerate}[label=\Alph*.,itemsep=0.1em]
  \item 寄存器
  \item L1 Cache
  \item L2 Cache
  \item DRAM（主存）
\end{enumerate}

\item 在链接过程中，关于强符号和弱符号的规则，说法正确的是（\hspace{1.2em}）。
\begin{enumerate}[label=\Alph*.,itemsep=0.1em]
  \item 未初始化的全局变量是强符号
  \item 初始化的全局变量是弱符号
  \item 函数名默认是强符号
  \item 局部静态变量在链接时参与全局符号解析，是弱符号
\end{enumerate}

\item 调用进程创建函数 \texttt{fork()} 之后，关于其返回值，下列说法正确的是（\hspace{1.2em}）。
\begin{enumerate}[label=\Alph*.,itemsep=0.1em]
  \item 在父进程和子进程中均返回 0
  \item 在父进程中返回 0，在子进程中返回父进程的 PID
  \item 在父进程中返回子进程的 PID，在子进程中返回 0
  \item 以上均不对
\end{enumerate}

\item 关于缺页（page fault），下列说法正确的是（\hspace{1.2em}）。
\begin{enumerate}[label=\Alph*.,itemsep=0.1em]
  \item 缺页属于异步异常（中断）
  \item 缺页属于故障（fault），处理完成后通常会重新执行引发缺页的指令
  \item 发生缺页时程序一定直接终止
  \item 缺页由用户程序直接处理，无需内核介入
\end{enumerate}

\item 关于虚拟内存，下列说法不正确的是（\hspace{1.2em}）。
\begin{enumerate}[label=\Alph*.,itemsep=0.1em]
  \item 每个进程拥有独立的虚拟地址空间
  \item DRAM 可看作对磁盘上虚拟地址空间内容的缓存
  \item 虚拟地址必须映射到物理地址的连续区域
  \item 虚拟内存有助于简化链接与加载、隔离进程
\end{enumerate}

\item （本题题干回忆不起来。）

\end{enumerate}

\section*{二、填空题}

\begin{enumerate}[resume]
\item 假设想用 16 个并行处理节点达到 10 的加速比，原计算程序中并行部分最少占比为 \blank。

\item 64 位系统中 \texttt{short} 数 $-11$ 的机器数二进制表示 \blanklong。

\item 链接器将可重定位目标文件生成可执行目标文件，主要经过 \blank 和 \blank 两个阶段。

\item \texttt{fork} 执行成功后返回 \blank 次；\texttt{execve} 调用成功则返回 \blank 次。

\item 虚拟内存系统借助 \blanklong 这一数据结构，将虚拟页映射到物理页。
\end{enumerate}

\section*{三、简答题}

\begin{enumerate}[resume]
\item 列举妨碍编译器进行程序优化的两个主要因素，并简述原理。

\vspace{3.5em}

\item 结合以下程序，解释局部性原理。
\begin{lstlisting}
int cal_array_sum(int *a, int n) {
    int sum = 0;
    for (int i = 0; i < n; i++)
        sum += a[i];
    return sum;
}
\end{lstlisting}

\vspace{2.5em}

\item 一函数的源程序以及进程图如下。请写出进程图中空白处的内容。
\begin{lstlisting}
void func() {
    printf("L0\n");
    if (fork() != 0) {
        printf("L1\n");
        if (fork() != 0) {
            printf("L2\n");
        }
    }
    printf("Bye\n");
}
\end{lstlisting}

\begin{center}
\includegraphics[width=0.92\textwidth]{process.png}
\end{center}

\noindent （1）\blanklong\qquad
（2）\blanklong\qquad
（3）\blanklong\\[0.4em]
（4）\blanklong\qquad
（5）\blanklong

\end{enumerate}

\section*{四、计算题}

\begin{enumerate}[resume]
\item 给出 C 语言变量 \texttt{float x = -0.2;} 在 IEEE 754 单精度浮点数格式下的十六进制表示。
要求写出推导过程（包括符号位、阶码、尾数的确定过程）。

\vspace{5em}

\item 某磁盘的平均旋转速率为 $10000\,\mathrm{RPM}$，平均寻道时间为 $5\,\mathrm{ms}$，每个扇区大小为 $512\,\mathrm{B}$，每个磁道平均扇区数为 $1000$。现有一个大小为 $10\,\mathrm{MB}$ 的文件连续存放在该磁盘上。请回答：
\begin{enumerate}
  \item 平均旋转延迟包括什么？
  \item 计算该磁盘的平均旋转延迟。
  \item 该 $10\,\mathrm{MB}$ 文件约占多少个扇区？（按 $1\,\mathrm{MB}=2^{20}\,\mathrm{B}$ 计算）
  \item 访问该文件的平均访问时间是多少？（写出计算过程）
\end{enumerate}

\vspace{4em}
\end{enumerate}

\section*{五、综合题}

\begin{enumerate}[resume]
\item 请写出 Y86-64 六阶段顺序执行（SEQ）处理器中 \texttt{call Dest} 指令每个阶段的微指令。

\begin{center}
\renewcommand{\arraystretch}{1.6}
\begin{tabular}{|p{2.2cm}|p{11cm}|}
\hline
\textbf{阶段} & \textbf{call Dest} \\
\hline
取指 & \\
\hline
译码 & \\
\hline
执行 & \\
\hline
访存 & \\
\hline
写回 & \\
\hline
更新 PC & \\
\hline
\end{tabular}
\end{center}

\item 假设某 CPU Cache 的参数为 $S=4$，$E=4$，$B=8$，若 $M=N=64$，请优化如下程序，并说明优化的方法（至少 CPU 与 Cache 各一种）。
\begin{lstlisting}
void trans(int M, int N, int A[M][N], int B[N][M]) {
    for (int i = 0; i < M; i++) {
        for (int j = 0; j < N; j++) {
            B[j][i] = A[i][j];
        }
    }
}
\end{lstlisting}

\vspace{3em}

\item 某计算机系统情况如下：
\begin{enumerate}[label=（\arabic*）,itemsep=0.1em]
  \item 内存按字节寻址，采用小端模式；
  \item 虚拟地址长度：16 位；
  \item 使用一级页表，页面大小为 256 字节；
  \item 物理地址长度为 16 位；
  \item TLB 为 4 路组相联，共 16 个条目；
  \item L1 d-Cache 为物理寻址、直接映射，行大小 4 字节，共 8 个组；
  \item TLB、Cache 的当前内容如表~\ref{tab:tlb}、表~\ref{tab:cache} 所示（数值均为十六进制）。
\end{enumerate}

\begin{table}[htbp]
\centering
\caption{TLB}
\label{tab:tlb}
\scriptsize
\renewcommand{\arraystretch}{1.15}
\begin{tabular}{|c|c|c|c|c|c|c|c|c|c|c|c|c|}
\hline
组 & 标记 & PPN & 有效 & 标记 & PPN & 有效 & 标记 & PPN & 有效 & 标记 & PPN & 有效 \\
\hline
0 & 03 & 1A & 1 & 12 & 2B & 0 & 07 & 3C & 1 & 1F & 0D & 0 \\
1 & 05 & 4E & 1 & 09 & 5F & 1 & 11 & 60 & 0 & 02 & 71 & 1 \\
2 & 0A & 82 & 0 & 14 & 93 & 1 & 08 & A4 & 1 & 16 & B5 & 0 \\
3 & 01 & C6 & 1 & 0C & D7 & 0 & 13 & E8 & 1 & 06 & F9 & 1 \\
\hline
\end{tabular}
\end{table}

\begin{table}[htbp]
\centering
\caption{L1 d-Cache}
\label{tab:cache}
\small
\renewcommand{\arraystretch}{1.15}
\begin{tabular}{|c|c|c|c|c|c|c|}
\hline
索引 & 标记 & 有效 & 字节0 & 字节1 & 字节2 & 字节3 \\
\hline
0 & 1A & 1 & 11 & 22 & 33 & 44 \\
1 & 4E & 1 & AA & BB & CC & DD \\
2 & 93 & 1 & 01 & 23 & 45 & 67 \\
3 & A4 & 1 & FE & DC & BA & 98 \\
4 & C6 & 1 & 10 & 32 & 54 & 76 \\
5 & E8 & 1 & 0F & 1E & 2D & 3C \\
6 & 71 & 1 & 55 & 66 & 77 & 88 \\
7 & F9 & 1 & 90 & A0 & B0 & C0 \\
\hline
\end{tabular}
\end{table}

CPU 要从虚拟地址 \texttt{0x05A3} 读取一个字（2 字节），用于后续计算。

请介绍：计算机和操作系统如何通过虚拟地址从内存中读取该值，并将该值返回给 CPU；同时结合表中 TLB（及 Cache）数据，求出最终读到的数值。
（叙述中应体现地址翻译、TLB、页表/缺页、物理寻址与 Cache/主存访问等关键环节。）

\vspace{6em}

\end{enumerate}

\end{document}
